\documentclass[11pt]{amsart}

\usepackage[T1]{fontenc}
\usepackage[utf8]{inputenc}
\usepackage{amsmath,amssymb,amsthm}
\usepackage{tikz}
\usepackage{booktabs}
\usepackage{array}
\usepackage[colorlinks=true,linkcolor=blue,citecolor=blue,urlcolor=blue]{hyperref}
\usepackage{microtype}

% Typewriter identifiers cannot hyphenate; allow interword stretch instead of
% overfull lines.
\setlength{\emergencystretch}{2.5em}

% ---------------------------------------------------------------------------
% Manually numbered theorem environments (numbering matches the source
% document and its internal cross-references exactly).
% ---------------------------------------------------------------------------
\theoremstyle{plain}
\newtheorem{innerthm}{Theorem}
\newenvironment{thmm}[2][]{%
  \renewcommand\theinnerthm{#2}%
  \ifx\relax#1\relax\begin{innerthm}\else\begin{innerthm}[#1]\fi
}{\end{innerthm}}

\newtheorem{innerprop}{Proposition}
\newenvironment{propm}[2][]{%
  \renewcommand\theinnerprop{#2}%
  \ifx\relax#1\relax\begin{innerprop}\else\begin{innerprop}[#1]\fi
}{\end{innerprop}}

\newtheorem{innercor}{Corollary}
\newenvironment{corm}[2][]{%
  \renewcommand\theinnercor{#2}%
  \ifx\relax#1\relax\begin{innercor}\else\begin{innercor}[#1]\fi
}{\end{innercor}}

\theoremstyle{definition}
\newtheorem{innerdefn}{Definition}
\newenvironment{defnm}[2][]{%
  \renewcommand\theinnerdefn{#2}%
  \ifx\relax#1\relax\begin{innerdefn}\else\begin{innerdefn}[#1]\fi
}{\end{innerdefn}}

\theoremstyle{remark}
\newtheorem{innerrem}{Remark}
\newenvironment{remm}[2][]{%
  \renewcommand\theinnerrem{#2}%
  \ifx\relax#1\relax\begin{innerrem}\else\begin{innerrem}[#1]\fi
}{\end{innerrem}}

% ---------------------------------------------------------------------------
% Notation macros
% ---------------------------------------------------------------------------
\newcommand{\ZZ}{\mathbb{Z}}
\newcommand{\Znn}[1]{Z_{#1}}
\newcommand{\himp}{\to}          % Heyting implication (relative pseudocomplement)
% Lean identifiers: typewriter face, breakable at underscores and slashes
% (url-package command: argument is read verbatim, no escaping of _ needed).
\DeclareUrlCommand\lean{\urlstyle{tt}%
  \def\UrlBreaks{\do\_\do\/\do\.\do\-}}

\begin{document}

\title[Unique ordinary elements, $\ZZ/12\ZZ$, and $n = p^2 q$]{The Unique
Ordinary Element of a One-Generated Heyting Algebra, the Subgroup Lattice of
$\ZZ/12\ZZ$, and a Characterization of $n = p^2 q$}

\author{Chris Brink}
\address{Independent researcher}
\email{chris@falsework.dev}
\urladdr{https://falsework.dev}

\date{July 2026. Preprint.}

\begin{abstract}
Call an element of a Heyting algebra \emph{ordinary} (following Citkin) if it
is neither regular ($\neg\neg a = a$) nor dense ($\neg a = \bot$). We collect,
with machine-checked proofs, a group of results centered on ordinary elements
and the Rieger--Nishimura lattice. (1)~\emph{Non-degeneracy:} relative to a
fixed element $a$, every non-bottom element of a Heyting algebra occupies
exactly one of four lattice regions determined by $a$, $\neg a$, and
$\neg\neg a$, and all four regions are simultaneously inhabited if and only if
$a$ is ordinary. (2)~\emph{Threshold:} any Heyting algebra containing an
ordinary element has at least six elements, and the six elements $\bot$, $g$,
$\neg g$, $\neg\neg g$, $g \vee \neg g$, $\top$ realize an order-embedding of
the six-element one-generated Heyting algebra $\Znn{6}$; the bound generalizes
an observation of Citkin from one-generated algebras with ordinary generators
to arbitrary Heyting algebras. (3)~\emph{Uniqueness:} in any one-generated
Heyting algebra of cardinality $\ge 6$ --- hence in every finite one-generated
algebra $\Znn{n}$ ($n \ge 6$) and in the free Heyting algebra on one generator
--- the generator is the unique ordinary element. The statement appears
without proof in a recent preprint of Citkin; the machine-checked proof given
here is, to our knowledge and his, the first written proof. En route we give a
machine-checked proof of Nishimura's one-variable normal-form theorem (every
element of the subalgebra generated by $g$ is a value of a Rieger--Nishimura
term in $g$). (4)~\emph{Converse failure:} uniqueness of the ordinary element
does not imply one-generation; the minimal counterexample has eight elements,
is unique at that cardinality, and is isomorphic to $\Znn{6}$ with a
three-element chain stacked above it. The construction class was independently
identified by Citkin. (5)~\emph{Arithmetic instance:} the subgroup lattice of
$\ZZ/12\ZZ$ (equivalently the divisor lattice of $12$) is one-generated as a
Heyting algebra --- by the two-element subgroup $\{0, 6\}$ --- and is
therefore isomorphic to $\Znn{6}$. Among all cyclic groups $\ZZ/n\ZZ$ with $n$
having at most two prime divisors, the subgroup lattice contains an ordinary
element iff $n$ is divisible by the square of a prime with a second prime
present, and contains a \emph{unique} ordinary element iff $n = p^2 q$; the
least such $n$ is $12$. In the pitch-class reading of $\ZZ/12\ZZ$, the unique
ordinary element is the tritone, and the four regions of (1) are inhabited by
the tritone, the augmented triad, the diminished seventh chord, and the
whole-tone scale. All principal results are kernel-checked in Lean~4 against
Mathlib4 with standard axioms only and no \lean{sorry}.
\end{abstract}

\maketitle

\section{Introduction}

The free Heyting algebra on one generator --- the Rieger--Nishimura lattice
(Rieger 1957; Nishimura 1960) --- is among the most completely understood
objects in intuitionistic logic. Citkin (2024) gives a modern algebraic
treatment, showing in particular that for each $n > 1$ there is exactly one
one-generated Heyting algebra $\Znn{n}$ of cardinality $n$, arising as the
finite quotients (truncations) of the free algebra. Citkin's paper isolates
the class of elements he calls \emph{ordinary} --- neither regular nor dense
--- and observes that $\Znn{2}$ through $\Znn{5}$ contain none, while the
generator of $\Znn{n}$ for $n \ge 6$ is ordinary.

This note collects a group of results that grew from a simple question: at
which elements $a$ of a Heyting algebra is the four-region lattice
decomposition
\begin{center}
below $a$ \;$\cdot$\; below $\neg a$ \;$\cdot$\; below $\neg\neg a$ but not
below $a$ \;$\cdot$\; meeting both $a$ and $\neg a$
\end{center}
\emph{non-degenerate}, in the sense that all four regions are simultaneously
inhabited? The answer (Theorem~3.2) is: exactly at the ordinary elements.
Ordinariness, introduced by Citkin as a structural condition on generators, is
thus equivalently the non-degeneracy condition for a natural partition of the
algebra --- the partition studied in a companion preprint in a topos-theoretic
setting (Brink 2026); the present note is purely lattice-theoretic and
self-contained.

From the non-degeneracy theorem, several consequences follow that we have not
found in the literature, and one that we subsequently found \emph{stated} but
not proved:

\begin{itemize}
\item \textbf{Theorem 4.1 (threshold).} Any Heyting algebra with an ordinary
element has at least six elements. This generalizes Citkin's
Proposition~4(c) (2024), which bounds one-generated algebras with ordinary
generators, to arbitrary Heyting algebras and arbitrary ordinary elements,
with an abstract proof.
\item \textbf{Theorem 4.2 (order-embedding).} For any ordinary $g$ in any
Heyting algebra $H$, the map $\bot \mapsto \bot$, $g \mapsto g$,
$\neg g \mapsto \neg g$, $\neg\neg g \mapsto \neg\neg g$,
$g \vee \neg g \mapsto g \vee \neg g$, $\top \mapsto \top$ is an
order-embedding of $\Znn{6}$ into $H$. The embedding is \emph{not} in general
a Heyting-subalgebra embedding (Remark~5.4 gives the counterexample in the
free algebra).
\item \textbf{Theorem 5.2 (uniqueness).} In any Heyting algebra one-generated
by an ordinary element $g$, the element $g$ is the \emph{unique} ordinary
element. The statement appears, without proof, in Citkin (2025, p.~13);
Citkin has confirmed (personal communication, 2026) that he knows of no
written proof. The machine-checked proof given here proceeds through a
machine-checked proof of Nishimura's one-variable normal-form theorem
(Theorem~5.1), which may be of independent interest as a formalization.
\item \textbf{Theorem 6.1 (converse failure).} Unique-ordinary does not imply
one-generated: there is an eight-element Heyting algebra $H_8$ with a unique
ordinary element that is not one-generated, and eight is the minimal such
cardinality. $H_8$ is $\Znn{6}$ with a three-element chain stacked above its
top; the construction class was independently identified by Citkin (personal
communication, 2026) from the statement alone.
\item \textbf{Theorem 7.1 (the arithmetic instance).} The subgroup lattice of
$\ZZ/12\ZZ$ is one-generated as a Heyting algebra by its two-element
subgroup, hence is isomorphic to $\Znn{6}$ by Citkin's uniqueness theorem;
its unique ordinary element is that subgroup. In the pitch-class reading this
subgroup is the tritone $\{0, 6\}$, and the four regions of the
non-degeneracy theorem are inhabited by four familiar objects of music theory
(Section~9).
\item \textbf{Theorem 8.1 (the $p^2 q$ law).} For $n$ with at most two prime
divisors, the subgroup lattice of $\ZZ/n\ZZ$ contains an ordinary element iff
$n$ is divisible by $p^2$ for some prime $p$ with a second prime present, and
contains a unique ordinary element iff $n = p^2 q$ ($p$, $q$ distinct
primes). The least $n$ of this form is $12$. Prime-power and squarefree
lattices never contain ordinary elements --- the former abstractly for all
chains, covering all $p^k$ at once.
\end{itemize}

The results are individually modest; their interest is in the convergence.
The smallest one-generated Heyting algebra with an ordinary element
($\Znn{6}$, forced by logic) and the smallest cyclic-group subgroup lattice
with a unique ordinary element ($n = 12$, forced by arithmetic) are the
\emph{same six-element algebra with the same distinguished element} --- and
that algebra is the pitch-class-set universe of twelve-tone equal temperament,
with the tritone as the distinguished element. Neither selection involves a
choice: the algebra is forced from two independent directions.

\subsection*{Formalization}
All principal results are formalized in Lean~4 against Mathlib4 and
kernel-checked; axiom audits report at most \lean{propext},
\lean{Classical.choice}, \lean{Quot.sound} (several results need only
\lean{propext}), and no proof contains \lean{sorry}. Section~10 gives the
theorem-by-theorem correspondence. The source is public
(\url{https://github.com/thefalsework/papers}, directory
\lean{lean/FalseWorkPapers/Examples/}).

\subsection*{Prior-art standing}
The prior-art situation was adjudicated by correspondence with Citkin in June
2026 (cited with permission as: A.~Citkin, personal communication, 2026-06-11)
and is stated precisely where each result appears. In summary: the threshold
facts for one-generated algebras are Citkin's (2024); the uniqueness statement
is Citkin's (2025, asserted without proof) and the written proof is apparently
new; the converse failure is new as a statement but was reconstructed by
Citkin in a few lines once posed, so we regard it as an observation, not an
advance; the $\ZZ/12\ZZ$ identification and the $p^2 q$ characterization did
not surface any prior art through the correspondence, though transmission
difficulties (Section~11) mean their novelty is unadjudicated rather than
established. Kurenniemi (2004) independently used divisor lattices as models
of musical structure, in just intonation and without the Heyting layer;
Section~9 records the relation.

\section{Preliminaries}

\subsection{Heyting algebras; regular, dense, and ordinary elements}

A \emph{Heyting algebra} $H$ is a bounded distributive lattice with a binary
operation $\himp$ (relative pseudocomplement) satisfying
$c \le a \himp b$ iff $c \wedge a \le b$. The \emph{pseudocomplement} is
$\neg a = a \himp \bot$. Standard identities used throughout:
$a \wedge \neg a = \bot$; $a \le \neg\neg a$; $\neg a = \neg\neg\neg a$;
$\neg\neg$ is a monotone idempotent operator.

Following Citkin (2024, \S2.1): an element $a \in H$ is \emph{regular} if
$\neg\neg a = a$, \emph{dense} if $\neg a = \bot$, and \textbf{ordinary} if it
is neither. Equivalently, $a$ is ordinary iff
\[
a \ne \bot \quad\text{and}\quad \neg a \ne \bot \quad\text{and}\quad
\neg\neg a \ne a.
\]
($a = \bot$ is regular, so the first clause is implied by the third and second
in context; we state all three because that is the form the non-degeneracy
criterion uses. The three clauses correspond to the classification of elements
as bottom-like, dense, or regular: an element is ordinary precisely when it
escapes all three.) In a Boolean algebra every element is regular, so Boolean
algebras have no ordinary elements.

\subsection{The Rieger--Nishimura lattice and the algebras $\Znn{n}$}

The free Heyting algebra $F(1)$ on one generator $p$ is the Rieger--Nishimura
lattice: its elements are $\bot$, $\top$, and the values of the
\emph{Nishimura terms} --- the ascending ladder generated from $p$ by the
recursion (in one standard indexing)
\[
g_0 = p, \quad g_1 = \neg p, \quad f_1 = p \vee \neg p, \quad
g_2 = \neg\neg p, \quad g_3 = \neg\neg p \himp p,
\]
\[
g_{n+4} = g_{n+3} \himp (g_n \vee g_{n+1}), \qquad
f_{n+2} = g_{n+2} \vee g_{n+1}.
\]

By Citkin (2024, Theorem~6), for each $n > 1$ there is exactly one
one-generated Heyting algebra of cardinality $n$, denoted $\Znn{n}$; the
$\Znn{n}$ are the finite quotients of $F(1)$. The first non-trivial ones:
\begin{itemize}
\item $\Znn{2} = \{\bot, \top\}$ (Boolean), $\Znn{3}$ = the three-element
chain, $\Znn{4}$, $\Znn{5}$ --- none of which contain an ordinary element
(Citkin 2024);
\item $\Znn{6}$, with elements $\bot$, $g$, $\neg g$, $\neg\neg g$,
$g \vee \neg g$, $\top$ and Hasse diagram
\end{itemize}

\begin{center}
\begin{tikzpicture}[scale=1.1, every node/.style={font=\small}]
  \node (bot) at (0.5, 0)   {$\bot$};
  \node (g)   at (0, 1)     {$g$};
  \node (ng)  at (1.5, 1)   {$\neg g$};
  \node (nng) at (-0.75, 2) {$\neg\neg g$};
  \node (gng) at (0.75, 2)  {$g \vee \neg g$};
  \node (top) at (0, 3)     {$\top$};
  \draw (bot) -- (g);
  \draw (bot) -- (ng);
  \draw (g)   -- (nng);
  \draw (g)   -- (gng);
  \draw (ng)  -- (gng);
  \draw (nng) -- (top);
  \draw (gng) -- (top);
\end{tikzpicture}
\end{center}

\noindent (the two incomparabilities are $\neg g \parallel \neg\neg g$ and
$\neg\neg g \parallel g \vee \neg g$; note $g \le \neg\neg g$ and both $g$,
$\neg g \le g \vee \neg g$).

We write ``$H$ is one-generated by $g$'' for: the smallest subset of $H$
containing $g$ and closed under $\wedge$, $\vee$, $\himp$, $\neg$, $\bot$,
$\top$ is $H$ itself.

\section{The four-region decomposition and the non-degeneracy theorem}

\begin{defnm}{3.1}
Fix a Heyting algebra $H$ and an element $a \in H$. For $x \in H$ with
$x \ne \bot$, say $x$ lies in region
\begin{itemize}
\item[\textbf{I}] if $x \le a$;
\item[\textbf{R}] if $x \le \neg a$;
\item[\textbf{E}] if $x \le \neg\neg a$ and $x \nleq a$;
\item[\textbf{D}] if $x \wedge a \ne \bot$ and $x \wedge \neg a \ne \bot$.
\end{itemize}
\end{defnm}

\begin{propm}[partition]{3.1}
For every $x \in H$ with $x \ne \bot$, exactly one of I, R, E, D holds.
\textup{(Lean: \lean{lattice_four_position_partition}.)}
\end{propm}

\begin{proof}
\emph{Disjointness.} I--R: $x \le a \wedge \neg a = \bot$. I--E: E requires
$x \nleq a$. I--D: $x \le a$ gives
$x \wedge \neg a \le a \wedge \neg a = \bot$. R--E and E--D:
$x \le \neg\neg a$ gives
$x \wedge \neg a \le \neg\neg a \wedge \neg a = \bot$, which contradicts
$x \le \neg a$ with $x \ne \bot$ (for R) and the second conjunct of D. R--D:
$x \le \neg a$ gives $x \wedge a = \bot$.
\emph{Exhaustiveness.} Case on whether $x \wedge a$ and $x \wedge \neg a$ are
$\bot$. If both are non-$\bot$, $x \in$ D. If $x \wedge \neg a = \bot$ then
$x \le \neg\neg a$, and $x \in$ I or $x \in$ E according as $x \le a$ or not.
If $x \wedge a = \bot$ then $x \le \neg a$, so $x \in$ R. If both are $\bot$
then $x \le \neg a \wedge \neg\neg a = \bot$, contradicting $x \ne \bot$.
\end{proof}

The decomposition is trivial for many choices of $a$: if $a = \bot$, regions
I, E, D are empty; if $a$ is dense, R is empty; if $a$ is regular, E is empty
($x \le \neg\neg a = a$ puts $x$ in I). The exact non-degeneracy condition is:

\begin{thmm}[non-degeneracy]{3.2}
In any Heyting algebra $H$ and for any $a \in H$, all four regions are
simultaneously inhabited if and only if $a$ is ordinary:
\[
a \ne \bot \;\wedge\; \neg a \ne \bot \;\wedge\; \neg\neg a \ne a.
\]
\textup{(Lean: \lean{allFourCellsInhabited_iff}.)}
\end{thmm}

\begin{proof}
($\Leftarrow$) The four witnesses are $a \in$ I, $\neg a \in$ R,
$\neg\neg a \in$ E (it is $\le \neg\neg a$; if $\neg\neg a \le a$ then
$\neg\neg a = a$, contradiction), and $a \vee \neg a \in$ D (its meets with
$a$ and $\neg a$ are $a \ne \bot$ and $\neg a \ne \bot$).
($\Rightarrow$) An inhabitant of I forces $a \ne \bot$; of R forces
$\neg a \ne \bot$; of E forces $\neg\neg a \ne a$ (an $x \le \neg\neg a$ with
$x \nleq a$ is impossible when $\neg\neg a = a$).
\end{proof}

Two immediate corollaries situate the condition:

\begin{corm}[Boolean collapse]{3.3}
No Boolean algebra contains an ordinary element; hence the decomposition is
never non-degenerate in a Boolean algebra. Moreover the double-negation map
$a \mapsto \neg\neg a$ of any Heyting algebra lands only on non-ordinary
(regular) elements, so the Boolean algebra of regular elements --- the
classical shadow of $H$ under Glivenko's translation --- never carries a
non-degenerate decomposition either.
\textup{(Lean: \lean{boolean_no_kernel}, \lean{glivenko_collapse}.)}
\end{corm}

\begin{corm}[dense-bottom lemma]{3.4}
If $H$ has a least non-zero element (an atom below every non-zero element),
then $H$ has no ordinary elements: every non-zero element is dense.
\textup{(Lean: \lean{no_ordinary_of_least_nonzero}.)}
\end{corm}

\section{The six-element threshold and the order-embedding}

\begin{thmm}[threshold]{4.1}
Let $H$ be a Heyting algebra containing an ordinary element $g$. Then the six
elements
\[
\bot, \quad g, \quad \neg g, \quad \neg\neg g, \quad g \vee \neg g, \quad \top
\]
are pairwise distinct; hence $|H| \ge 6$ (in particular any finite Heyting
algebra with a non-degenerate four-region decomposition has at least six
elements).
\textup{(Lean: \lean{ordinary_card_ge_six}, \lean{allFourCells_card_ge_six}.)}
\end{thmm}

\begin{proof}[Proof sketch]
Most distinctions are immediate from ordinariness ($g \ne \bot$,
$\neg g \ne \bot$, $\neg\neg g \ne g$) and the basic identities
($g \wedge \neg g = \bot$ separates $g$ from $\neg g$ since both are
non-$\bot$; $\neg\neg g \wedge \neg g = \bot$ similarly). The load-bearing
case is $g \vee \neg g \ne \top$: if $g \vee \neg g = \top$ then, by
distributivity,
$\neg\neg g = \neg\neg g \wedge (g \vee \neg g)
= (\neg\neg g \wedge g) \vee (\neg\neg g \wedge \neg g) = g \vee \bot = g$,
contradicting ordinariness.
\end{proof}

\subsection*{Relation to prior art}
For \emph{one-generated} algebras with ordinary \emph{generators} this is
Citkin's Proposition~4(c) (2024), and the non-existence of ordinary elements
in $\Znn{2}$--$\Znn{5}$ is observed in the same paper. Theorem~4.1 removes
both restrictions: any ordinary element, any Heyting algebra, abstract proof
with no enumeration. The $\Znn{5}$ case then follows as a corollary rather
than a check (the formalization derives it both ways and confirms agreement).

\begin{thmm}[order-embedding]{4.2}
With $H$ and $g$ as above, the assignment
\[
\bot \mapsto \bot, \quad g \mapsto g, \quad \neg g \mapsto \neg g, \quad
\neg\neg g \mapsto \neg\neg g, \quad
(g \vee \neg g) \mapsto g \vee \neg g, \quad \top \mapsto \top
\]
defines an order-embedding $\Znn{6} \hookrightarrow H$: the six images realize
exactly the Hasse structure of $\Znn{6}$, including the two incomparabilities
$\neg g \parallel \neg\neg g$ and $\neg\neg g \parallel g \vee \neg g$ (each
of the $18$ non-relations is forced by ordinariness). The axiom audit for this
theorem reports \lean{propext} only.
\textup{(Lean: \lean{div12OrderEmbedding}.)}
\end{thmm}

\subsection*{Scope guard}
Theorem~4.2 is an \emph{order}-embedding, not a Heyting-subalgebra embedding:
the six-element image is in general not closed under $\himp$. Remark~5.4
exhibits the failure in $F(1)$ itself. (A note on the Lean name:
\lean{div12OrderEmbedding} is stated for an arbitrary Heyting algebra $H$ and
arbitrary ordinary $g$ --- the name refers to the \emph{source} of the
embedding, $\mathrm{Div}12 \cong \Znn{6}$, not to an instance restriction.)
The subalgebra-level statement that \emph{is} true is Citkin's
Proposition~3.1 (2025): any Heyting algebra with an ordinary element contains
a subalgebra isomorphic to some $\Znn{n}$, $n \ge 6$ --- the subalgebra
generated by the ordinary element, which by Theorem~5.1 below is a quotient of
the Rieger--Nishimura ladder, not necessarily the six-element truncation.

\section{Nishimura normal form and the uniqueness of the ordinary element}

The engine of the uniqueness proof is the classical one-variable normal-form
theorem, which we formalize in full:

\begin{thmm}[Nishimura normal form, machine-checked]{5.1}
Let $H$ be a Heyting algebra and $g \in H$. Every element of the Heyting
subalgebra generated by $g$ is $\top$ or a value at $g$ of a Nishimura term.
\textup{(Lean: \lean{generatedBy_isLadderValue}.)}
\end{thmm}

\begin{proof}[Proof structure (as formalized)]
The set of ladder values is shown closed under each Heyting operation in turn:
the complement table ($\neg x_n$ is again a ladder value, with
$\neg x_n = \bot$ from index~4 up); the join and meet tables (comparable pairs
are absorbed; the incomparable diagonals are the recursion identities
$x_m \vee x_{m+1} = x_{m+3}$ and $x_m \vee x_{m+2} = x_{m+5}$ for $m$ odd,
dually for meets); and the implication table, whose intricacy is the real
content. The decisive simplification is that implications into
\emph{odd-indexed} targets are terminal ---
$x_m \himp x_{2i+3} = x_{2i+3}$ whenever $x_{2i+1} \le x_m$, by currying ---
so the only genuinely recursive family is implication into even targets, which
climbs by two with a comparability ceiling; a lexicographic measure closes the
induction, and the Lean kernel certifies termination. Structural induction
over the generated subalgebra then gives the theorem. Axioms: \lean{propext},
\lean{Quot.sound}; no \lean{Classical.choice}.
\end{proof}

The theorem itself is classical (Nishimura 1960; algebraic proof in Citkin
2024). Its formal value here is that it converts ``one-generated'' into an
induction principle, discharging what was previously a citation from the
following result:

\begin{thmm}[uniqueness of the ordinary element]{5.2}
Let $H$ be a Heyting algebra one-generated by $g$, and suppose $g$ is
ordinary. Then $g$ is the unique ordinary element of $H$. Consequently every
$\Znn{n}$ with $n \ge 6$, and the free Heyting algebra $F(1)$ itself, has its
generator as unique ordinary element --- equivalently (Theorem~3.2), as the
unique element at which the four-region decomposition is non-degenerate.
\textup{(Lean: \lean{nishimura_ordinary_unique},
\lean{nishimura_kernel_unique_of_generated}.)}
\end{thmm}

\begin{proof}[Proof idea]
By Theorem~5.1 every element is a ladder value; the ladder is then analyzed
directly. Every rung from index~4 up lies above $g \vee \neg g$ and is
therefore dense (its pseudocomplement is $\bot$); the rungs $\neg g$ and
$\neg\neg g$ are regular; $\bot$ and $\top$ are regular and dense
respectively. Only the generator's rung survives. Verified additionally by
unconditional exhaustive computation at $n = 6, 7, 8$. Axioms:
\lean{propext}, \lean{Quot.sound}.
\end{proof}

\subsection*{Prior-art standing (adjudicated)}
The statement of Theorem~5.2 appears in print in Citkin (2025, p.~13): ``the
algebra $Z_n$ contains a single ordinary element that generates the entire
algebra'' --- asserted without proof, feeding his Proposition~3.1. In
correspondence, Citkin stated he is ``not aware of any known proofs'' and
sketched the folklore argument specialists would give: the free algebra
$F(1)$ has a single ordinary element; the finite $\Znn{n}$ are its quotients;
regularity and density are preserved under Heyting quotients; hence uniqueness
passes down. That argument presupposes the single-ordinary-element property of
$F(1)$ itself, which is exactly what the ladder analysis above establishes;
the two routes are consistent. The standing of Theorem~5.2 is therefore:
\emph{statement published by Citkin as an unproved observation; the
machine-checked proof presented here apparently the first written proof.}

\begin{corm}[unique generator]{5.3}
A one-generated Heyting algebra of cardinality $\ge 6$ has a unique generator.
(Generators are ordinary by Citkin's Proposition~4(c); the ordinary element is
unique by Theorem~5.2. The general statement was communicated by Citkin; the
instance at $\Znn{6}$ is kernel-checked directly as
\lean{Div12.generator_unique}, by trapping each non-generator inside an
explicit proper subalgebra.)
\end{corm}

\begin{remm}[the subalgebra upgrade of Theorem~4.2 is false in general]{5.4}
In $F(1)$, with $g$ the free generator, $\neg\neg g \himp g$ is the fifth rung
of the ladder --- a genuinely new element outside the six-element set
$\{\bot, g, \neg g, \neg\neg g, g \vee \neg g, \top\}$. So the order-embedding
of Theorem~4.2 does not preserve $\himp$ in general; it is a Heyting embedding
exactly when the ladder above index~4 collapses into the six-set
(equivalently, when the subalgebra generated by $g$ has six elements). This is
a kernel-checked counterexample, recorded to prevent the natural overreading
of Theorem~4.2.
\end{remm}

\section{The converse fails: $H_8$}

Theorem~5.2 says one-generated (with ordinary generator) $\Rightarrow$ unique
ordinary element. The converse is false, and the minimal counterexample is
small:

\begin{thmm}[converse failure]{6.1}
There is an eight-element Heyting algebra $H_8$ with a unique ordinary element
$a$ such that $H_8$ is not one-generated by $a$ (nor by any element). Eight is
minimal: every Heyting algebra of cardinality $\le 7$ with a unique ordinary
element is one-generated --- concretely, is $\Znn{6}$ or $\Znn{7}$, since
$\Znn{2}$--$\Znn{5}$ contain no ordinary element. $H_8$ is unique at
cardinality eight up to isomorphism.
\textup{(Lean: \lean{unique_ordinary_converse_false}.)}
\end{thmm}

\begin{proof}[Construction]
$H_8$ is $\Znn{6}$ with a three-element chain stacked above its top (the top
of $\Znn{6}$ identified with the bottom of the chain): $6 + 3 - 1 = 8$
elements. The two new elements above the old top are dense, so no new ordinary
elements are created and the unique ordinary element of $\Znn{6}$ survives;
but the term ladder of $a$ stabilizes after seven values and never reaches the
eighth element, so $a$ does not generate (machine-checked: the term values of
$a$ stay inside a proper subalgebra --- which, per Citkin's Proposition~3.1,
is precisely a $\Znn{7}$ core; the isomorphism of the core with $\Znn{7}$ is
itself kernel-checked, \lean{h8_ladder_core}).
\end{proof}

\subsection*{Status of the minimality claim}
The non-generation and the $\Znn{7}$-core anatomy are kernel-checked. The
\emph{minimality} (no counterexample below eight elements, uniqueness at
eight) rests on an exhaustive enumeration of finite Heyting algebras of
cardinality $\le 12$ --- via Birkhoff duality, as downset lattices of finite
posets, with per-size counts cross-checked against OEIS A006982 --- performed
by script, not formalized in Lean. We flag the different evidential grade
explicitly.

\subsection*{Independent confirmation}
Posed the bare statement, Citkin reconstructed the construction class from
first principles without seeing the witness: glue a one-generated six-element
algebra to a chain of height $\ge 3$; new chain elements are dense; the unique
ordinary element survives; the result is not one-generated; smallest case
$\Znn{6} + \Znn{3}$, eight elements (personal communication, 2026-06-11). His
reconstruction matches $H_8$ exactly. He regarded the fact as readily
derivable once posed; we therefore present Theorem~6.1 as an observation with
a machine-checked witness, not as a field advance.

\begin{propm}[dense-filter salvage]{6.2}
In any Heyting algebra with a unique ordinary element $a$, the dense elements
are exactly the filter above $a \vee \neg a$.
\textup{(Lean: \lean{unique_ordinary_dense_iff}.)}
\end{propm}

\begin{proof}[Proof sketch (following the formalization)]
Let $x$ be dense. Meeting with a dense element does not change complements:
$\neg(b \wedge x) = \neg b$. Hence $\neg(a \wedge x) = \neg a \ne \bot$ and
$\neg\neg(a \wedge x) = \neg\neg a$, which cannot equal $a \wedge x$ (that
would force $\neg\neg a \le a$, collapsing $a$'s ordinariness); so
$a \wedge x$ is ordinary, hence $a \wedge x = a$ by uniqueness, i.e.
$a \le x$. For $\neg a$: the element $\neg a \wedge x$ has complement
$\neg\neg a \ne \bot$ (since $a \ne \bot$), so if it were not regular it would
be ordinary, hence equal to $a$ by uniqueness --- forcing $a \le \neg a$ and
so $a = \bot$, a contradiction. So $\neg a \wedge x$ is regular:
$\neg a \wedge x = \neg\neg(\neg a \wedge x) = \neg\neg\neg a = \neg a$, i.e.
$\neg a \le x$. Hence $a \vee \neg a \le x$. Conversely anything above
$a \vee \neg a$ is dense, since
$\neg(a \vee \neg a) = \neg a \wedge \neg\neg a = \bot$.
\end{proof}

The statement was independently confirmed by Citkin (personal communication,
2026), by an argument of the same shape. The least dense element is thus
$a \vee \neg a$ --- the fourth Nishimura term.

The pair $(\Znn{8}, H_8)$ is instructive: same cardinality, both with a unique
ordinary element, only one of them one-generated. Uniqueness of the ordinary
element does not certify that an algebra is a truncation of the free algebra;
the two properties of $\Znn{6}$ established in the next section --- unique
ordinary element \emph{and} one-generation --- are independent facts requiring
separate proof.

\section{The subgroup lattice of $\ZZ/12\ZZ$ is $\Znn{6}$}

Let $\mathrm{Div}12$ denote the divisor lattice of $12$ --- elements
$\{1, 2, 3, 4, 6, 12\}$, order = divisibility, meet = $\gcd$, join =
$\mathrm{lcm}$ --- which is isomorphic to the subgroup lattice of the cyclic
group $\ZZ/12\ZZ$ under $d \mapsto$ (the unique subgroup of order $d$).
$\mathrm{Div}12$ is a finite distributive lattice, hence a Heyting algebra.

\begin{thmm}[the arithmetic instance]{7.1}
$\mathrm{Div}12$ is one-generated as a Heyting algebra by the divisor $2$
(equivalently, the two-element subgroup $\{0, 6\} \le \ZZ/12\ZZ$). Explicitly,
with $g = 2$:
\[
\neg g = 3, \quad \neg\neg g = 4, \quad g \vee \neg g = 6, \quad
g \wedge \neg g = 1 = \bot, \quad \neg\neg g \vee (g \vee \neg g) = 12 = \top.
\]
\textup{(Lean: \lean{Div12.one_generated_by_tritone},
\lean{Div12.tritone_generates}.)}
\end{thmm}

Hence, by Citkin's uniqueness theorem (2024, Theorem~6),
$\mathbf{Div12 \cong \Znn{6}}$, with the divisor $2$ as the (unique, by
Corollary~5.3) free generator. By Theorem~5.2 the divisor $2$ is also the
unique ordinary element of $\mathrm{Div}12$; the formalization derives this
three independent ways (exhaustive check, the uniqueness law, and through the
generation machinery), and the three agree.

The proof is a finite computation, but the statement has content: it
identifies a lattice arising in nature (subgroups of a cyclic group;
pitch-class sets closed under transposition, in the reading of Section~9) with
the smallest non-degenerate truncation of the free Heyting algebra on one
generator, and its distinguished subgroup with the free generator. Neither
side of the identification was constructed to meet the other.

\section{Divisor lattices of $\ZZ/n\ZZ$: the $p^2 q$ characterization}

Which cyclic groups have subgroup lattices containing ordinary elements? The
subgroup lattice of $\ZZ/n\ZZ$ is the divisor lattice of $n$, which is a
product of chains --- one chain per prime divisor, of length (exponent~+~1).
(This structure theorem is standard and is cited, not re-proved, for general
$n$; the instance at $n = 12$, $\mathrm{Div}12 \cong C_3 \times C_2$ as
Heyting algebras, is kernel-checked with $\himp$ and $\neg$ preserved,
\lean{div12OrderIsoChains}.)

\begin{thmm}[the $p^2 q$ law]{8.1}
\
\begin{enumerate}
\item A totally ordered Heyting algebra contains no ordinary element (every
non-$\bot$ element has pseudocomplement $\bot$, i.e.\ is dense). Hence no
prime-power $n = p^k$ yields an ordinary element --- abstractly, for all $k$.
\item In a product of two non-trivial bounded chains, the ordinary elements
are exactly the pairs with one coordinate $\bot$ and the other strictly
between $\bot$ and $\top$ of its chain.
\item For $n = p^a q^b$ ($a, b \ge 1$, $p \ne q$ prime): the divisor lattice
of $n$ contains an ordinary element iff $\max(a, b) \ge 2$. In particular
squarefree $n = pq$ ($6, 10, 15, \dots$) yield none.
\item The ordinary element is \emph{unique} iff $\{a, b\} = \{2, 1\}$, i.e.\
iff $n = p^2 q$. The least such $n$ is $12$.
\item Boundary behavior, kernel-checked at instances:
$n = 24 = 2^3 \cdot 3$ has two ordinary elements,
$n = 36 = 2^2 \cdot 3^2$ has two,
$n = 60 = 2^2 \cdot 3 \cdot 5$ has three.
\end{enumerate}
\textup{(Lean: \lean{total_no_kernel}, \lean{chainProd_kernel_exists_iff},
\lean{chainProd_kernel_unique_iff}, \lean{why_twelve}.)}
\end{thmm}

Under the explicit isomorphism $\mathrm{Div}12 \cong C_3 \times C_2$, the
unique ordinary element of the $n = 12$ lattice is carried to $(1, 0)$ --- and
back in $\mathrm{Div}12$ it is the divisor $2$, the generator of Theorem~7.1.

\begin{remm}[the count formula, two primes only]{8.2}
Statement~(2) yields an exact count: the divisor lattice of $n = p^a q^b$ has
precisely $(a - 1) + (b - 1)$ ordinary elements --- the strictly internal
elements of one chain paired with $\bot$ of the other. This gives (3)--(5) at
once for two primes: existence iff the count is positive iff
$\max(a, b) \ge 2$; uniqueness iff the count is $1$ iff
$\{a, b\} = \{2, 1\}$; and $24 = 2^3 \cdot 3 \to 2$,
$36 = 2^2 \cdot 3^2 \to 2$. The formula does \emph{not} extend coordinatewise
to three or more primes: for $60 = 2^2 \cdot 3 \cdot 5$ the naive sum
$\sum_i (a_i - 1)$ gives $1$, while the kernel-checked count is $3$
(\lean{sixty_kernel_three}) --- in a product of three chains an ordinary
element needs only \emph{some} coordinate at $\bot$ alongside a strictly
internal one, not all remaining coordinates at $\bot$, so $(m, 0, 0)$,
$(m, 1, 0)$, $(m, 0, 1)$ all qualify. This is why the $\ge 3$-prime case is
left open below rather than claimed by extension.
\end{remm}

\subsection*{Scope statements, fixed precisely}
(i)~The characterization in (3)--(4) is proved abstractly for all exponents
but only for $n$ with \textbf{at most two prime divisors}; for three or more
primes we have instance checks only ($n = 60$), and the general $k$-prime
statement is left open (the coordinatewise computation appears routine but is
not formalized). (ii)~$18 = 2 \cdot 3^2$ and $20 = 2^2 \cdot 5$ are also of
the form $p^2 q$ and also carry unique ordinary elements --- necessarily the
\emph{same} six-element lattice. Twelve is the \textbf{least} $n$ of the form,
not the only one. (iii)~The identification of $n$-tone equal temperament with
$\ZZ/n\ZZ$ (next section) is a modeling convention, not mathematics.

\subsection*{The two-sided forcing}
Theorems~5.2 and~8.1 select the same object from independent directions. On
the logic side, $\Znn{6}$ is the \emph{first} one-generated Heyting algebra
whose four-region decomposition is non-degenerate (Theorem~4.1: nothing below
six elements qualifies), with the free generator as the forced distinguished
element. On the arithmetic side, $n = 12$ is the \emph{first} cyclic-group
order whose subgroup lattice carries an ordinary element at all, and --- $12$
being of the uniqueness shape $p^2 q$ --- carries exactly one. The two firsts
coincide: $\Znn{6} \cong \mathrm{Div}12 \cong C_3 \times C_2$, free generator
= divisor $2$ = $(1, 0)$.

\section{A remark on the musical reading}

This section is interpretive and is flagged as such; nothing in
Sections~3--8 depends on it.

Identify $\ZZ/12\ZZ$ with the pitch classes of twelve-tone equal temperament.
The subgroups are then the transpositionally symmetric pitch-class sets:
$\{0\}$ (trivial), $\{0, 6\}$ (the \textbf{tritone} dyad), $\{0, 4, 8\}$ (the
\textbf{augmented triad}), $\{0, 3, 6, 9\}$ (the \textbf{diminished seventh}
chord), $\{0, 2, 4, 6, 8, 10\}$ (the \textbf{whole-tone} scale), and
$\ZZ/12\ZZ$ (the chromatic aggregate). Under Theorem~7.1 the unique ordinary
element --- the free generator of the lattice qua $\Znn{6}$ --- is the
tritone. The four regions of the non-degeneracy theorem (Theorem~3.2) at the
tritone are inhabited by:

\begin{center}
\begin{tabular}{lll}
\toprule
region & witness & pitch-class object \\
\midrule
I (below $g$) & $g$ & tritone $\{0, 6\}$ \\
R (below $\neg g$) & $\neg g$ & augmented triad $\{0, 4, 8\}$ \\
E (below $\neg\neg g$, not below $g$) & $\neg\neg g$ &
  diminished seventh $\{0, 3, 6, 9\}$ \\
D (meets both) & $g \vee \neg g$ & whole-tone scale $\{0, 2, 4, 6, 8, 10\}$ \\
\bottomrule
\end{tabular}
\end{center}

That the diminished seventh --- the pivot chord of tonal ambiguity in
common-practice harmony --- occupies the double-negation residue
$\neg\neg g \setminus g$, the region that exists only in non-Boolean logic, is
at minimum a striking coincidence of structure; we do not claim more for it
here. A companion series of framework papers (Brink 2026, and references
therein) develops an interpretive program on top of this instance; the present
note is confined to the mathematics.

\subsection*{Prior art in the musical direction}
Kurenniemi (2004) independently proposed divisor lattices as models of musical
structure, working in 5-limit just intonation with the divisor lattices of
large numbers (e.g.\ 8640) under $\gcd$/$\mathrm{lcm}$, and observing the
near-coincidence $2^{19} \approx 3^{12}$ (the Pythagorean comma). His
treatment stops at the distributive lattice and its geometric embedding: no
Heyting structure, no relative pseudocomplement, no one-generation, no
ordinary elements. The results of Sections~7--8 do not appear there; we cite
it as independent precedent for the divisor-lattice formalism in music, with
the observation that his machinery aimed at $n = 12$ rather than
just-intonation integers would meet the present construction.

\section{Formalization}

All principal results are formalized in Lean~4 against Mathlib4 (toolchain
\lean{v4.30.0-rc2} era) and kernel-checked. No proof contains \lean{sorry}; no
proof uses \lean{native_decide}. Axiom audits (\texttt{\#print axioms}) report:

\begin{center}
\scriptsize
\renewcommand{\arraystretch}{1.25}
\begin{tabular}{@{}>{\raggedright\arraybackslash}p{7.5em}
                 >{\raggedright\arraybackslash}p{11.5em}
                 >{\raggedright\arraybackslash}p{10.5em}
                 >{\raggedright\arraybackslash}p{5.5em}@{}}
\toprule
Result & Lean theorem & File (\lean{Examples/}) & Axioms \\
\midrule
Partition (Prop.~3.1) & \lean{lattice_four_position_partition} &
  \lean{Lattice/FourPositionLattice.lean}\textsuperscript{$\dagger$} &
  all three \\
Non-degeneracy (Thm~3.2) & \lean{allFourCellsInhabited_iff} &
  \lean{NishimuraKernelLaw.lean} & \lean{propext} \\
Boolean/Glivenko (Cor.~3.3) & \lean{boolean_no_kernel},
  \lean{glivenko_collapse} &
  \lean{GlivenkoCollapse.lean} & \lean{propext} \\
Dense-bottom (Cor.~3.4) & \lean{no_ordinary_of_least_nonzero} &
  \lean{LadderCore.lean} & all three \\
Threshold (Thm~4.1) & \lean{ordinary_card_ge_six},
  \lean{allFourCells_card_ge_six} &
  \lean{LadderCore.lean} & all three \\
Order-embedding (Thm~4.2) & \lean{div12OrderEmbedding} &
  \lean{LadderCore.lean} & \lean{propext} \\
Normal form (Thm~5.1) & \lean{generatedBy_isLadderValue} &
  \lean{NishimuraNormalForm.lean} &
  \lean{propext}, \lean{Quot.sound} \\
Uniqueness (Thm~5.2) & \lean{nishimura_ordinary_unique},
  \lean{nishimura_kernel_unique_of_generated} &
  \lean{NishimuraKernelLaw.lean}, \lean{NishimuraNormalForm.lean} &
  \lean{propext}, \lean{Quot.sound} \\
Unique generator at $\Znn{6}$ (Cor.~5.3) & \lean{Div12.generator_unique} &
  \lean{LadderCore.lean} & \lean{propext}, \lean{Quot.sound} \\
Subalgebra failure (Rem.~5.4) & via \lean{nishimuraTerm_himp_diagonal} &
  \lean{NishimuraNormalForm.lean} &
  \lean{propext}, \lean{Quot.sound} \\
$H_8$ (Thm~6.1) & \lean{unique_ordinary_converse_false},
  \lean{h8_ladder_core} &
  \lean{UniqueOrdinaryConverse.lean}, \lean{LadderCore.lean} &
  \lean{propext}, \lean{Quot.sound} \\
Dense filter (Prop.~6.2) & \lean{unique_ordinary_dense_iff} &
  \lean{UniqueOrdinaryConverse.lean} &
  \lean{propext}, \lean{Quot.sound} \\
$\mathrm{Div}12 \cong \Znn{6}$ (Thm~7.1) &
  \lean{Div12.one_generated_by_tritone},
  \lean{Div12.tritone_generates}, \lean{Div12.kernel_unique}
  (+ \lean{_via_law}, \lean{_via_subalgebra}) &
  \lean{NishimuraTruncations.lean}, \lean{NishimuraNormalForm.lean} &
  \lean{propext}, \lean{Quot.sound} \\
$p^2 q$ law (Thm~8.1) & \lean{total_no_kernel},
  \lean{chainProd_kernel_exists_iff},
  \lean{chainProd_kernel_unique_iff}, \lean{why_twelve},
  \lean{div12OrderIsoChains} &
  \lean{WhyTwelve.lean} & all three \\
\bottomrule
\end{tabular}
\end{center}

\noindent ``All three'' abbreviates \lean{propext}, \lean{Classical.choice},
\lean{Quot.sound}. \textsuperscript{$\dagger$}The partition theorem lives
under \lean{lean/FalseWorkPapers/Lattice/}, not \lean{Examples/}.

Truncation computations at $\Znn{5}$--$\Znn{8}$ are exhaustive \lean{decide}
checks over concrete decidable instances certified one-generated by exhibiting
every element as a term value. Two items rest partly outside the kernel and
are flagged where stated: the minimality of $H_8$ (script-level enumeration,
Section~6) and the $\ge 3$-prime case of Theorem~8.1 (instances only). The
general ``divisor lattice = product of chains'' structure theorem and Citkin's
Theorem~6 (uniqueness of $\Znn{n}$) are cited classical inputs.

Source: \url{https://github.com/thefalsework/papers}, in the directory
\lean{lean/FalseWorkPapers/Examples/}; audit lines are collected in
\lean{Examples/HeytingTypeInstance.lean}.

\subsection*{Methodological note}
Several proofs were developed against the published Rieger--Nishimura
structure rather than from finite models, for a reason worth recording: the
finite truncations $\Znn{n}$ are quotients of $F(1)$, so an identity read off
a finite model may be a collapse artifact --- true in $\Znn{n}$, false in
$F(1)$ --- and would compile. (Concretely: $\Znn{8}$ reports
$x_5 \vee x_6 = \top$, but in $F(1)$ the join is the eighth rung.) Finite
models were used as sound \emph{refuters} only, within their faithful window;
every load-bearing table entry is proved abstractly.

\section{Prior art, correspondence, and acknowledgements}

The prior-art adjudication was conducted by correspondence with Alex Citkin
(June 2026, three exchanges; cited with his permission as personal
communication, 2026-06-11). His 2024 paper supplies the $\Znn{n}$
classification, the ordinary-element notion, the $\Znn{2}$--$\Znn{5}$
observation, and Proposition~4(c); his December 2025 preprint (Citkin 2025)
states the uniqueness of the ordinary element without proof (p.~13) and proves
Proposition~3.1 (subalgebra containment of a $\Znn{n}$, $n \ge 6$, in any
algebra with an ordinary element). In correspondence he confirmed knowing no
written proof of the uniqueness statement, supplied the folklore quotient
argument, independently reconstructed the $H_8$ construction class and
confirmed its eight-element minimality, communicated the unique-generator fact
(Corollary~5.3), and supplied an independent proof of Proposition~6.2. The
identification $\mathrm{Div}12 \cong \Znn{6}$ and the $p^2 q$ characterization
were posed to him but, owing to a terminological ambiguity in the posing
(``divisor lattice of 12'' read as a twelve-element algebra), were not
adjudicated; their novelty standing is accordingly \emph{no prior art
surfaced, not established}. We record this rather than claim novelty.

The question that motivated the non-degeneracy theorem arose in a companion
program (Brink 2026) concerning a four-fold partition of subobject lattices in
topoi; the present note isolates everything that is pure lattice theory. Erkki
Kurenniemi's unpublished 2004 notebook was located after the fact and is cited
as independent precedent for divisor-lattice models of musical structure
(Section~9).

\subsection*{AI disclosure}
This paper and its formalization were produced with AI assistance (Anthropic
Claude) under the correction architecture documented in a companion
methodology paper (Brink 2026b): machine verification for every principal
claim, explicit evidential grading of unformalized steps, pre-registered
experiments for the truncation computations, and expert correspondence for
prior-art adjudication.

\begin{thebibliography}{99}

\bibitem{BBdJ2008}
Bezhanishvili, G., Bezhanishvili, N., \& de Jongh, D. (2008).
The Kuznetsov--Ger\v{c}iu and Rieger--Nishimura logics: The boundaries of the
finite model property.
\emph{Logic and Logical Philosophy}, 17.
(For the Nishimura polynomial recursion.)

\bibitem{Birkhoff1937}
Birkhoff, G. (1937).
Rings of sets.
\emph{Duke Mathematical Journal}, 3(3), 443--454.

\bibitem{Brink2026}
Brink, C. (2026).
A four-position partition of morphisms in elementary topoi with distinction
structure.
Preprint, \url{https://github.com/thefalsework/papers}.

\bibitem{Brink2026b}
Brink, C. (2026b).
Epistemic dependency as structural condition.
Preprint, \url{https://github.com/thefalsework/papers}.

\bibitem{Citkin2024}
Citkin, A. (2024).
An algebraic proof of the Nishimura theorem.
\emph{Logics}, 2(4), 148--157.

\bibitem{Citkin2025}
Citkin, A. (2025).
Hereditarily structurally complete superintuitionistic logics and primitive
varieties of Heyting algebras.
arXiv:2512.05633.

\bibitem{Citkin2026}
Citkin, A. (2026).
Personal communication, 2026-06-11. Cited with permission.

\bibitem{dJC2006}
de Jongh, D., \& Chagrova, L. (2006).
ILLC Prepublication PP-2006-25, University of Amsterdam.
(Nishimura polynomials, Definition 37.)

\bibitem{Kurenniemi2004}
Kurenniemi, E. (2004).
\emph{Chords, scales, and divisor lattices.}
Unpublished Mathematica notebook, DIMI.

\bibitem{Mathlib}
Mathlib Community (2026).
Mathlib4.
\url{https://github.com/leanprover-community/mathlib4}.

\bibitem{Nishimura1960}
Nishimura, I. (1960).
On formulas of one variable in intuitionistic propositional calculus.
\emph{Journal of Symbolic Logic}, 25.

\bibitem{OEIS}
OEIS Foundation (n.d.).
A006982: Number of unlabeled distributive lattices with $n$ elements.
\url{https://oeis.org/A006982}.

\bibitem{Rieger1957}
Rieger, L. (1957).
A remark on the so-called free closure algebras.
\emph{Czechoslovak Mathematical Journal}, 7.

\end{thebibliography}

\end{document}
